Here is a corrected proof that makes the common-domain/compatibility and \(\sigma\)-finiteness hypotheses explicit.

The sum-space form of Riesz--Thorin is the appropriate version here: it assumes \(\sigma\)-finite measure spaces and a single linear operator
\[
T:L^{p_0}(X)+L^{p_1}(X)\to L^{q_0}(Y)+L^{q_1}(Y),
\]
so the two endpoint actions are compatible because they are restrictions of the same \(T\). This is the form stated, for example, in Tao’s notes on interpolation of \(L^p\) spaces. ([metaphor.ethz.ch](https://metaphor.ethz.ch/x/2017/hs/401-3370-67L/sc/Tao.pdf))

**Riesz--Thorin interpolation theorem.** Let \((X,\mu)\) and \((Y,\nu)\) be \(\sigma\)-finite measure spaces, and let
\[
1\le p_0,p_1,q_0,q_1\le \infty.
\]
For \(0<\theta<1\), define \(p_\theta,q_\theta\) by
\[
\frac1{p_\theta}
=
\frac{1-\theta}{p_0}
+
\frac{\theta}{p_1},
\qquad
\frac1{q_\theta}
=
\frac{1-\theta}{q_0}
+
\frac{\theta}{q_1},
\]
with the convention \(1/\infty=0\). Suppose
\[
T:L^{p_0}(X)+L^{p_1}(X)\to L^{q_0}(Y)+L^{q_1}(Y)
\]
is linear, and suppose that for \(i=0,1\),
\[
T\bigl(L^{p_i}(X)\bigr)\subseteq L^{q_i}(Y),
\qquad
\|Tf\|_{q_i}\le M_i\|f\|_{p_i}
\quad\text{for all }f\in L^{p_i}(X).
\]
Then
\[
T\bigl(L^{p_\theta}(X)\bigr)\subseteq L^{q_\theta}(Y),
\]
and
\[
\|T\|_{L^{p_\theta}\to L^{q_\theta}}
\le
M_0^{1-\theta}M_1^\theta .
\]

Now let
\[
\Omega=\{1,\dots,m\}
\]
with counting measure \(\mu\). Then
\[
\mu(\Omega)=m<\infty,
\]
so \((\Omega,\mu)\) is \(\sigma\)-finite. Indeed,
\[
\Omega=\bigcup_{j=1}^m \{j\},
\qquad
\mu(\{j\})=1<\infty.
\]
For every \(1\le p\le \infty\),
\[
L^p(\Omega,\mu)=\ell_p^m
\]
as a vector space, with the usual \(\ell_p\)-norm.

Define a single linear operator
\[
T:\mathbb C^m\to \mathbb C^m,
\qquad
Tx=Hx.
\]
Since
\[
L^2(\Omega,\mu)=L^\infty(\Omega,\mu)=\mathbb C^m
\]
as sets, we have
\[
L^2(\Omega,\mu)+L^\infty(\Omega,\mu)=\mathbb C^m.
\]
Thus \(T\) is one linear map on the common domain
\[
L^2+L^\infty,
\]
and its endpoint actions on \(L^2\) and \(L^\infty\) are restrictions of the same operator. This verifies the compatibility hypothesis.

We are given
\[
\|H\|_{2\to 2}=\sqrt m,
\qquad
\|H\|_{\infty\to\infty}=m.
\]
Hence the Riesz--Thorin hypotheses hold with
\[
p_0=q_0=2,
\qquad
p_1=q_1=\infty,
\qquad
M_0=\sqrt m,
\qquad
M_1=m.
\]

First, for \(r=2\),
\[
\|H\|_{2\to 2}
=
\sqrt m
=
m^{1/2}
=
m^{1-1/2},
\]
so the desired estimate holds.

Now let \(r>2\). Choose
\[
\theta=1-\frac{2}{r}.
\]
Since \(r>2\), we have \(0<\theta<1\). Also,
\[
\frac1{p_\theta}
=
\frac{1-\theta}{2}
+
\frac{\theta}{\infty}
=
\frac{1-\theta}{2}
=
\frac1r,
\]
and similarly
\[
\frac1{q_\theta}
=
\frac1r.
\]
Therefore
\[
p_\theta=q_\theta=r.
\]
Applying Riesz--Thorin gives
\[
\|H\|_{r\to r}
\le
\|H\|_{2\to 2}^{1-\theta}
\|H\|_{\infty\to\infty}^{\theta}
=
(\sqrt m)^{1-\theta}m^\theta.
\]
Using
\[
1-\theta=\frac{2}{r},
\qquad
\theta=1-\frac{2}{r},
\]
we obtain
\[
(\sqrt m)^{1-\theta}m^\theta
=
(\sqrt m)^{2/r}m^{1-2/r}
=
m^{1/r}m^{1-2/r}
=
m^{1-1/r}.
\]
Thus
\[
\boxed{\|H\|_{r\to r}\le m^{1-1/r}}
\]
for every real \(r\ge 2\).

If one also allows the endpoint \(r=\infty\), the same formula gives
\[
\|H\|_{\infty\to\infty}=m=m^{1-1/\infty},
\]
using the convention \(1/\infty=0\).